% =====================================================================
%  The Physiology of Response:
%  When an Interaction Is Relevant to an Agent, and What a
%  Relevant Response Costs
%
%  Self-contained. States its own axioms; re-derives every primitive
%  it uses; cites only standard, externally published mathematics.
% =====================================================================
\documentclass[11pt,twocolumn,a4paper]{article}

\usepackage[utf8]{inputenc}
\usepackage[T1]{fontenc}
\usepackage{lmodern}
\usepackage[margin=1in]{geometry}
\usepackage{amsmath,amssymb,amsthm,mathtools}
\usepackage{thmtools}
\usepackage{enumitem}
\usepackage{microtype}
\usepackage{booktabs}
\usepackage{graphicx}
\usepackage{caption}
\usepackage{hyperref}
\usepackage{cleveref}
\usepackage{xcolor}
\usepackage{tikz}
\usetikzlibrary{arrows.meta,positioning,calc}

\hypersetup{
  colorlinks=true, linkcolor=blue!50!black,
  citecolor=green!50!black, urlcolor=blue!60!black
}

% ---- theorem environments ----
\declaretheoremstyle[
  spaceabove=6pt, spacebelow=6pt,
  headfont=\bfseries, notefont=\mdseries, notebraces={(}{)},
  bodyfont=\itshape, postheadspace=1em
]{thmstyle}
\declaretheoremstyle[
  spaceabove=6pt, spacebelow=6pt,
  headfont=\bfseries, notefont=\mdseries, notebraces={(}{)},
  bodyfont=\normalfont, postheadspace=1em
]{defstyle}

\declaretheorem[style=thmstyle, name=Theorem, numberwithin=section]{theorem}
\declaretheorem[style=thmstyle, name=Proposition, sibling=theorem]{proposition}
\declaretheorem[style=thmstyle, name=Lemma, sibling=theorem]{lemma}
\declaretheorem[style=thmstyle, name=Corollary, sibling=theorem]{corollary}
\declaretheorem[style=defstyle, name=Definition, sibling=theorem]{definition}
\declaretheorem[style=defstyle, name=Axiom, sibling=theorem]{axiom}
\declaretheorem[style=defstyle, name=Remark, sibling=theorem]{remark}
\declaretheorem[style=defstyle, name=Example, sibling=theorem]{example}
\declaretheorem[style=defstyle, name=Principle, sibling=theorem]{principle}

\crefname{axiom}{Axiom}{Axioms}
\Crefname{axiom}{Axiom}{Axioms}

% ---- notation ----
\newcommand{\R}{\mathbb{R}}
\newcommand{\N}{\mathbb{N}}
\newcommand{\Agent}{\mathsf{A}}
\newcommand{\Beh}{\mathcal{B}}       % behaviour / disposition space
\newcommand{\Purp}{x^{\star}}        % purpose point
\newcommand{\drive}{\Phi}            % drive potential
\newcommand{\lam}{\lambda}
\newcommand{\Int}{\iota}             % an interaction
\newcommand{\resp}{r}                % a response
\newcommand{\Resp}{\mathcal{R}}      % response repertoire
\newcommand{\floorc}{\beta}          % response floor
\newcommand{\att}{\alpha}            % attention budget
\newcommand{\cost}{c}                % response cost
\newcommand{\gain}{g}                % purpose gain
\newcommand{\coh}{\kappa}            % coherence margin
\newcommand{\rel}{\mathrm{rel}}
\newcommand{\abs}[1]{\lvert#1\rvert}
\newcommand{\norm}[1]{\lVert#1\rVert}
\newcommand{\dist}{d}
\newcommand{\inner}[2]{\langle #1,#2\rangle}

\title{\textbf{The Physiology of Response:\\
When an Interaction Is Relevant to an Agent,\\
and What a Relevant Response Costs}}

\author{Kundai Farai Sachikonye\\[2pt]
\small AIMe Registry for Artificial Intelligence\\
\small \texttt{kundai.sachikonye@bitspark.com}}
\date{\today}

\begin{document}
\maketitle

\begin{abstract}
\noindent
An agent embedded in a busy world is addressed constantly, and most of
what reaches it does not matter to it. We ask what it is for an incoming
\emph{interaction} to be \emph{relevant} to an agent, and what a relevant
response \emph{costs}, and we answer with mathematics a reader can follow
without any other work. We take an agent to have two standing quantities,
which we define from scratch: a \emph{purpose}, the attractor of its own
motivation, and a \emph{coherence}, the internal consistency it must not
break. A candidate response to an interaction is admissible only if it is
\emph{two-factor relevant}: it must (a) advance the agent toward its
purpose and (b) preserve the agent's coherence. From four axioms---bounded
attention, a positive floor on the cost of any response, purpose-directed
drive, and coherence as a closed consistency condition---we prove:
\textbf{(T1)} relevance is the conjunction of a purpose gain and a
coherence margin, and neither alone suffices; \textbf{(T2)} an agent has a
bounded attention budget, so it can commit to at most finitely many
responses per unit time and must decline the rest---inattention is
forced, not a defect; \textbf{(T3)} the agent's optimal policy is a
threshold rule: respond exactly to interactions whose purpose gain exceeds
a price set by the attention budget, and ignore the rest, which reproduces
the lived phenomenology of an agent that is present but preoccupied;
\textbf{(T4)} coherence is testable locally and in advance---an
incoherent response is detectable from the signs of its internal effects
alone, before it is committed, so an agent can refuse a purpose-advancing
response that would break it; \textbf{(T5)} response is floored and
irreversible: every committed response deposits a positive, unrepealable
change in the agent, so an agent cannot un-receive what it has responded
to, and successive encounters are never with the same agent;
\textbf{(T6)} construction and commitment are mutually exclusive at each
instant---an agent forming new distinctions cannot simultaneously commit
to an action---so responsiveness and reflection alternate, which is a
second, independent source of the appearance of distraction. We give the
admissible-response construction, the threshold policy, the coherence
test, and a numerical study on explicit agents. The reading of the formal
apparatus as attention, preoccupation, and the felt autonomy of a
character is marked as such; no theorem depends on it.
\end{abstract}

\vspace{0.5em}
\noindent\textbf{Keywords:} relevance; attention budget; threshold policy;
coherence; holonomy; irreversibility; construction--commitment
alternation; non-player character; agent response.

% =====================================================================
\section{Introduction}
\label{sec:intro}
% =====================================================================

\subsection{The question}

An agent in a shared world receives far more than it can act on. A passer
speaks to it; an event occurs nearby; an object enters its field. Most of
this is nothing to the agent---not because the agent fails to perceive it,
but because responding to it would not serve what the agent is doing and
is not worth the agent's limited capacity. The mark of a convincing
embedded agent is precisely this selectivity: it is present, but it is not
at anyone's disposal. It answers what matters to it and lets the rest
pass, sometimes with a token acknowledgement, sometimes with none.

This paper asks what makes an incoming interaction \emph{relevant} to an
agent, and what a relevant response \emph{costs}. We want a definition
sharp enough to compute a policy from---given an interaction, should the
agent respond, and if so how---and self-contained: nothing but this paper
and standard mathematics.

\subsection{The two-factor thesis}

We argue that relevance has two factors and needs both. A response to an
interaction is worth making only if it \emph{advances the agent toward its
purpose} and \emph{preserves the agent's coherence}. The first factor is
self-interest: a response that does not move the agent toward what it is
trying to do is, to the agent, idle. The second is integrity: a response
that would advance the purpose but break the agent's internal
consistency---make it act against itself---is inadmissible. Relevance is
the conjunction. An interaction to which no response passes both tests is
irrelevant to the agent, and the correct behaviour is to decline it.

From this, together with the plain fact that the agent's capacity to
respond is bounded, the phenomenology of an embedded character follows:
selective attention, declined interactions, token acknowledgements, and
the alternation between engaging and withdrawing.

\subsection{Relation to the companion account of purpose}

This paper is self-contained; it defines every quantity it uses. It does,
however, take from a companion study the picture of an agent as having a
standing \emph{identity} and a \emph{purpose} that attracts its conduct.
We restate exactly what we need---purpose as the attractor of a drive, and
coherence as a closed consistency condition---as \Cref{sec:standing}
axioms and definitions, proved or posited here, so that the present paper
stands alone. A reader who has not seen the companion account loses
nothing.

\subsection{Organisation}

\Cref{sec:standing} fixes the agent's standing quantities---purpose and
coherence---and the interaction and response. \Cref{sec:relevance} proves
the two-factor relevance theorem (T1). \Cref{sec:attention} proves the
bounded-attention and threshold-policy results (T2, T3).
\Cref{sec:coherence} proves that coherence is testable in advance (T4).
\Cref{sec:irreversible} proves response is floored and irreversible (T5).
\Cref{sec:phases} proves construction--commitment alternation (T6).
\Cref{sec:validation} reports a numerical study. \Cref{sec:discussion}
discusses scope.

% =====================================================================
\section{The agent's standing quantities, the interaction, the response}
\label{sec:standing}
% =====================================================================

We fix the objects. The agent has a space of dispositions, a purpose it is
drawn to, and a coherence it must keep. An interaction arrives; a response
is a change the agent may make to itself in reply.

\subsection{Dispositions, purpose, drive}

\begin{definition}[Disposition space]
\label{def:disp}
The \emph{disposition space} $\Beh$ of an agent is a compact convex subset
of $\R^n$; a point $x\in\Beh$ is the agent's current disposition (its
tendencies to act). Distance on $\Beh$ is Euclidean, $\dist(x,y)=
\norm{x-y}$.
\end{definition}

\begin{axiom}[Purpose-directed drive]
\label{ax:drive}
The agent carries a \emph{drive potential} $\drive:\Beh\to\R$,
continuously differentiable and $\lam$-strongly convex for some $\lam>0$:
$\drive(y)\ge\drive(x)+\nabla\drive(x)^\top(y-x)+\tfrac{\lam}{2}
\norm{y-x}^2$ for all $x,y$. Its behaviour, absent interaction, follows
the descent $\dot x=-\nabla\drive(x)$.
\end{axiom}

\begin{definition}[Purpose]
\label{def:purpose}
The \emph{purpose} of the agent is the unique minimiser
$\Purp=\arg\min_{x\in\Beh}\drive(x)$; it is the attractor of the drive.
For a disposition $x$, the \emph{residual to purpose} is
$D(x)=\drive(x)-\drive(\Purp)\ge0$, the drive-height still to be
descended.
\end{definition}

\begin{lemma}[Purpose is well-defined and attracting]
\label{lem:purpose}
Under \Cref{ax:drive}, $\Purp$ exists and is unique, and along the drive
$\norm{x(t)-\Purp}\le e^{-\lam t}\norm{x(0)-\Purp}$; equivalently
$D(x(t))\le e^{-2\lam t}D(x(0))$.
\end{lemma}

\begin{proof}
A continuous function on a compact set attains its minimum; strong
convexity makes the minimiser unique. For attraction, let
$V=\tfrac12\norm{x-\Purp}^2$; then $\dot V=-(x-\Purp)^\top\nabla\drive(x)
\le-\lam\norm{x-\Purp}^2=-2\lam V$ using strong-convexity monotonicity and
$\nabla\drive(\Purp)=0$, so $V(t)\le e^{-2\lam t}V(0)$ by Grönwall. The
$D$-form follows since $D$ is comparable to $\tfrac{\lam}{2}\norm{x-\Purp}^2$
from below and to $\tfrac{L}{2}\norm{x-\Purp}^2$ from above for the
gradient's Lipschitz constant $L$.
\end{proof}

\subsection{Coherence}

An agent is coherent when its dispositions do not contradict one another:
when the ways it is disposed to act, taken around any internal loop of
dependence, agree. We formalise this as the vanishing of a
\emph{holonomy}---the net drift accumulated on going around a cycle---which
is the standard way to say ``all routes agree.''

\begin{definition}[Consistency graph; holonomy; coherence]
\label{def:coherence}
The agent's \emph{consistency graph} $H$ has a vertex for each disposition
component and a directed edge $i\to j$ carrying a \emph{transport}
$\theta_{ij}\in\R^d$ recording the shift the agent's structure induces in
$j$ per unit of $i$. The \emph{holonomy} of a directed cycle
$c=(i_0,i_1,\dots,i_0)$ is $\eta(c)=\sum_\ell\theta_{i_\ell i_{\ell+1}}$.
The agent is \emph{coherent} if $\eta(c)=0$ for every cycle $c$; the
\emph{coherence margin} of a disposition $x$ is
$\coh(x)=-\max_c\norm{\eta_x(c)}$, the negative of the worst cycle
disagreement at $x$ (so $\coh(x)=0$ at perfect coherence and $\coh(x)<0$
when some cycle fails to close).
\end{definition}

\begin{axiom}[Coherence is a closed condition]
\label{ax:closed}
The set of coherent dispositions $\{x\in\Beh:\coh(x)=0\}$ is closed and
nonempty, and $\coh:\Beh\to(-\infty,0]$ is continuous. (An agent has at
least one coherent disposition and coherence does not jump.)
\end{axiom}

\begin{remark}[Why holonomy is the right notion]
\label{rem:holonomy}
``All routes agree'' is exactly ``every cycle sums to zero.'' If two chains
of internal dependence lead from a disposition to another and disagree,
their difference is a cycle with nonzero holonomy; conversely a nonzero
cycle is two routes that disagree. Coherence as vanishing holonomy is
therefore not a special assumption but the general form of internal
consistency, and---crucially for \Cref{sec:coherence}---it is checkable
from the \emph{signs and sums of local transports}, without any global
solve.
\end{remark}

\subsection{Interaction and response}

\begin{definition}[Interaction]
\label{def:interaction}
An \emph{interaction} $\Int$ is an external event presented to the agent:
formally a map $\Int:\Beh\to\Beh$ that the agent \emph{may} apply to its
disposition (the change the interaction would induce if the agent takes it
up), together with a tag identifying it. An interaction is \emph{received}
whether or not it is responded to; it is \emph{responded to} only if the
agent applies a response to it.
\end{definition}

\begin{definition}[Response; repertoire; cost]
\label{def:response}
A \emph{response} to $\Int$ is a map $\resp:\Beh\to\Beh$ the agent applies
to its own disposition in reply, drawn from a finite \emph{repertoire}
$\Resp$. Applying $\resp$ moves the disposition from $x$ to
$\resp\!\cdot x$. Each response has a \emph{cost} $\cost(\resp)>0$,
the capacity it consumes.
\end{definition}

\begin{axiom}[Response floor]
\label{ax:floor}
No response is free: there is $\floorc>0$ with $\cost(\resp)\ge\floorc$ for
every $\resp\in\Resp$. Even the smallest genuine response---a token
acknowledgement---costs at least $\floorc$.
\end{axiom}

\begin{axiom}[Bounded attention]
\label{ax:attention}
The agent has a finite \emph{attention budget} $\att>0$ per unit time: the
total cost of the responses it commits in any unit interval cannot exceed
$\att$. Formally, if $R$ is the set of responses committed in a unit
interval, $\sum_{\resp\in R}\cost(\resp)\le\att$.
\end{axiom}

% =====================================================================
\section{Two-factor relevance}
\label{sec:relevance}
% =====================================================================

We now define when an interaction is relevant and prove that relevance is
the conjunction of a purpose gain and a coherence margin, with neither
sufficient alone.

\begin{definition}[Purpose gain]
\label{def:gain}
The \emph{purpose gain} of a response $\resp$ to $\Int$ at disposition $x$
is the residual it removes,
\[
\gain(\resp;x)\ =\ D(x)\ -\ D(\resp\!\cdot x)\ =\
\drive(x)-\drive(\resp\!\cdot x),
\]
the drop in drive-height. The response is \emph{purpose-advancing} at $x$
if $\gain(\resp;x)>0$.
\end{definition}

\begin{definition}[Coherence-preserving]
\label{def:cohpres}
A response $\resp$ is \emph{coherence-preserving} at $x$ if it does not
worsen the coherence margin: $\coh(\resp\!\cdot x)\ge\coh(x)$, i.e. the
agent is no less coherent after responding than before. It is
\emph{coherence-breaking} if $\coh(\resp\!\cdot x)<\coh(x)$.
\end{definition}

\begin{definition}[Two-factor relevance]
\label{def:relevance}
An interaction $\Int$ is \emph{relevant} to the agent at disposition $x$
if there exists a response $\resp\in\Resp$ that is both purpose-advancing
and coherence-preserving at $x$:
\[
\exists\,\resp\in\Resp:\quad
\gain(\resp;x)>0\ \ \text{and}\ \ \coh(\resp\!\cdot x)\ge\coh(x).
\]
Such a response is \emph{admissible}. An interaction with no admissible
response is \emph{irrelevant} to the agent.
\end{definition}

\begin{theorem}[Relevance is a conjunction; neither factor alone suffices]
\label{thm:conjunction}
Under \Cref{ax:drive,ax:closed}:
\begin{enumerate}[label=\textup{(\roman*)},leftmargin=2.2em]
\item An interaction is relevant iff its admissible set is nonempty, i.e.
iff some response clears \emph{both} the purpose-gain and
coherence-margin bars.
\item Purpose gain alone does not make an interaction relevant: there
exist interactions whose every purpose-advancing response is
coherence-breaking, hence inadmissible.
\item Coherence preservation alone does not make an interaction relevant:
a response that preserves coherence but has $\gain\le0$ is inadmissible.
\end{enumerate}
\end{theorem}

\begin{proof}
(i) is \Cref{def:relevance} restated: relevance is nonemptiness of the
admissible set, and admissibility is the conjunction of the two
conditions. (ii) Construct $\Beh=[0,1]^2$, a strongly convex $\drive$ with
minimiser at a corner, and a repertoire whose only $\gain>0$ response
carries the disposition across a locus where some cycle's holonomy jumps
from $0$ to a positive norm, so $\coh$ strictly decreases; then every
purpose-advancing response is coherence-breaking, the admissible set is
empty, and the interaction is irrelevant despite offering purpose gain.
(Continuity of $\coh$, \Cref{ax:closed}, makes the decrease well-defined.)
(iii) If $\gain(\resp;x)\le0$ then $\resp$ is not purpose-advancing, so it
fails the first conjunct and is inadmissible regardless of its coherence
effect.
\end{proof}

\begin{remark}[The peripheral player, derived]
\label{rem:peripheral}
\Cref{thm:conjunction} already yields the central phenomenological fact.
An interaction from another party (a player addressing the agent) is
relevant only if some response to it both advances the agent's own purpose
and keeps the agent coherent. When the player's address does neither---as
is typical when the agent is busy with its own pursuit---the interaction
is simply irrelevant, and the agent's correct behaviour is to not respond,
or to respond only with the cheapest coherence-neutral token. The player
is peripheral by construction: relevance is measured against the agent's
purpose, not the player's. No special ``distraction'' mechanism is needed;
peripherality is what relevance-to-self means.
\end{remark}

% =====================================================================
\section{Bounded attention and the threshold policy}
\label{sec:attention}
% =====================================================================

Relevance says which interactions \emph{could} be worth a response.
Bounded attention says the agent cannot take up all of them, and forces a
selection. We show the selection is a threshold rule and that declining
interactions is not a failure but the optimal use of a finite budget.

\begin{theorem}[Inattention is forced]
\label{thm:inattention}
Under \Cref{ax:floor,ax:attention}, in any unit interval the agent commits
at most $\lfloor\att/\floorc\rfloor$ responses. If more than
$\lfloor\att/\floorc\rfloor$ relevant interactions arrive in the interval,
the agent must leave at least one relevant interaction without response.
Declining interactions is therefore forced by finiteness, not a defect of
the agent.
\end{theorem}

\begin{proof}
Each committed response costs $\ge\floorc$ (\Cref{ax:floor}) and the total
committed cost is $\le\att$ (\Cref{ax:attention}); hence the number of
committed responses $N$ satisfies $N\floorc\le\att$, i.e.
$N\le\lfloor\att/\floorc\rfloor$. If the number of relevant interactions
exceeds this bound, not all can be committed, so at least one relevant
interaction goes without response.
\end{proof}

\begin{definition}[Response policy; value]
\label{def:policy}
A \emph{response policy} selects, from the relevant interactions arriving
in a unit interval, a subset $R$ to respond to, subject to
$\sum_{\resp\in R}\cost(\resp)\le\att$. Its \emph{value} is the total
purpose gain secured, $\sum_{\resp\in R}\gain(\resp;x)$ (each interaction
answered by its best admissible response).
\end{definition}

\begin{theorem}[Optimal policy is a threshold on gain density]
\label{thm:threshold}
The value-maximising policy under the attention budget is a threshold
rule: order relevant interactions by \emph{gain density}
$\gain(\resp;x)/\cost(\resp)$ and commit them in decreasing order until the
budget $\att$ is exhausted; equivalently, there is a price
$p^\star\ge0$ such that the agent responds to an interaction iff its best
admissible response has gain density $\ge p^\star$, and ignores the rest.
The price $p^\star$ rises as the budget $\att$ falls or as the arrival of
relevant interactions grows.
\end{theorem}

\begin{proof}
This is the fractional/continuous knapsack structure: maximising
$\sum_{\resp\in R}\gain(\resp)$ subject to $\sum_{\resp\in R}\cost(\resp)
\le\att$ over selections $R$. Sorting by gain density and taking greedily
until the budget is spent is optimal for the continuous relaxation, and
the Lagrangian dual assigns a single multiplier $p^\star$---the shadow
price of attention---such that an item is selected iff its gain density
exceeds $p^\star$ (items exactly at $p^\star$ fill the residual budget).
Tightening $\att$ or adding higher-density competitors raises the binding
$p^\star$. The integral (all-or-nothing) version is approximated by the
same threshold up to one boundary item, and coincides with it when
responses are divisible into acknowledgement-sized units of cost
$\floorc$.
\end{proof}

\begin{remark}[Interpretation: present but preoccupied]
\label{rem:preoccupied}
\Cref{thm:threshold} is the policy of an agent that is present but
preoccupied. It answers the interactions that most advance its purpose per
unit of attention and lets the rest pass, with a price $p^\star$ that
climbs exactly when the agent is busy (low spare $\att$) or besieged (many
competing relevant interactions). A player addressing such an agent during
a high-$p^\star$ moment receives, correctly, a thin response or none---not
because the agent is broken, but because the address fell below the
agent's current price of attention. Lowering $p^\star$ (an idle agent)
makes the same agent forthcoming. The single scalar $p^\star$ is thus a
readable dial for how reachable a character is at a given moment.
\end{remark}

\begin{corollary}[Token responses are the boundary case]
\label{cor:token}
The cheapest admissible response---cost $\floorc$, coherence-neutral,
minimal positive gain---is committed exactly when the budget has residual
capacity below the cost of any higher-gain response but at least $\floorc$.
This is the formal ``one sec'' / nod: the agent spends its last unit of
attention on a floor-cost acknowledgement rather than a substantive reply.
\end{corollary}

\begin{proof}
By \Cref{thm:threshold} the agent fills the budget greedily; when the
residual budget lies in $[\floorc,\,\min_{\text{substantive}}\cost)$, the
only affordable admissible response is a floor-cost one, which it takes if
its gain density clears $p^\star$.
\end{proof}

% =====================================================================
\section{Coherence is testable in advance}
\label{sec:coherence}
% =====================================================================

The second relevance factor would be useless if the agent could not tell,
\emph{before} committing, whether a response will break it. We show
coherence is decidable locally and in advance, from the signs and sums of
the response's internal effects, without a global solve.

\begin{theorem}[Coherence is a local, ordinal, pre-commitment test]
\label{thm:cohtest}
A response $\resp$ is coherence-preserving at $x$ iff, for every cycle $c$
in the consistency graph, the post-response holonomy has norm no larger
than the pre-response holonomy:
$\norm{\eta_{\resp\cdot x}(c)}\le\norm{\eta_x(c)}$. This is decidable from
the transports $\theta_{ij}$ alone---their signs and directed sums around
cycles---without computing the agent's full future behaviour, and it can
be evaluated before $\resp$ is committed.
\end{theorem}

\begin{proof}
By \Cref{def:coherence}, $\coh(y)=-\max_c\norm{\eta_y(c)}$, so
$\coh(\resp\!\cdot x)\ge\coh(x)$ iff $\max_c\norm{\eta_{\resp\cdot x}(c)}
\le\max_c\norm{\eta_x(c)}$, which holds iff no cycle's holonomy norm
increases past the prior maximum; the stated per-cycle inequality is
sufficient and, taken over the maximising cycle, necessary. Each
$\eta_y(c)=\sum_\ell\theta_{i_\ell i_{\ell+1}}$ is a finite directed sum of
local transports, computable from the graph's edge data at $y$ and at
$\resp\!\cdot x$; a cycle basis (of size $\abs{E}-\abs{V}+1$ for a
connected consistency graph) suffices, so the test is polynomial in the
graph and uses only local transport data, evaluable on the hypothetical
post-response state without committing to it.
\end{proof}

\begin{corollary}[An agent can refuse a tempting but destructive response]
\label{cor:refuse}
An agent can decline a response that advances its purpose ($\gain>0$) but
fails the coherence test, before committing it. Thus the agent is not
forced to pursue purpose at the cost of self-consistency: two-factor
relevance (\Cref{def:relevance}) is enforceable in advance, because its
second factor is a pre-commitment, local, ordinal check.
\end{corollary}

\begin{proof}
By \Cref{thm:cohtest} the coherence factor is evaluable before commitment;
if it fails while $\gain>0$, the response is purpose-advancing but
coherence-breaking, hence inadmissible by \Cref{def:relevance}, and the
agent declines it on the basis of the pre-commitment test alone.
\end{proof}

\begin{principle}[Integrity over opportunity]
\label{prin:integrity}
Because coherence is testable in advance, an agent modelled this way will
pass up purpose-advancing interactions that would make it incoherent---it
will act against an immediate gain to remain itself. This is the formal
counterpart of a character with integrity: not every advantageous
interaction is taken, and which are refused is decidable and inspectable.
\end{principle}

% =====================================================================
\section{Response is floored and irreversible}
\label{sec:irreversible}
% =====================================================================

We now show that responding changes the agent permanently: every committed
response deposits at least the floor of change and cannot be undone, so an
agent cannot un-receive what it responded to, and no two encounters are
with the same agent.

\begin{definition}[Committed history; agent state]
\label{def:history}
The agent carries a \emph{committed count} $m\in\N$, incremented by one
each time it commits a response. Its \emph{state} is the pair
$(x,m)$ of current disposition and committed count. Two states are equal
iff both components agree.
\end{definition}

\begin{theorem}[Floored, irreversible response]
\label{thm:irreversible}
Every committed response deposits a change of size $\ge\floorc$ in the
agent's expended attention and increments the committed count by one.
Consequently:
\begin{enumerate}[label=\textup{(\roman*)},leftmargin=2.2em]
\item the committed count is strictly increasing, so for encounters at
counts $m<m'$ the agent's states differ, $(x,m)\neq(x',m')$, even if
$x=x'$;
\item there is no response that restores a prior state: an ``undo'' is
itself a committed response, raising the count further; hence a revisited
disposition is never a revisited state;
\item an agent cannot un-receive an interaction it has responded to: the
response's deposit and the count increment persist.
\end{enumerate}
\end{theorem}

\begin{proof}
Committing a response costs $\ge\floorc$ (\Cref{ax:floor}) and, by
\Cref{def:history}, increments $m$ by one; so counts strictly increase.
(i) If $m\neq m'$ the second components differ, so the states differ
regardless of dispositions. (ii) An operation intended to restore an
earlier state is a further committed response, producing count $m+1\neq m$;
it yields a new state, not the old one. (iii) The deposit (expended
attention) and the count increment are not removed by any subsequent
response, since every response only adds to the count; so the fact of
having responded persists.
\end{proof}

\begin{remark}[Each encounter is with a new agent]
\label{rem:new-agent}
\Cref{thm:irreversible} says that an agent a player meets twice is, at the
second meeting, in a strictly later state than at the first---its committed
history has grown---even if its disposition looks the same. The agent that
returns is not the agent that left; resemblance is not identity of state.
This is what gives an embedded character the feel of a continuing life
rather than a resetting script: responding is a one-way deposit, and the
agent accumulates.
\end{remark}

% =====================================================================
\section{Construction and commitment alternate}
\label{sec:phases}
% =====================================================================

Finally we show a second, independent source of the appearance of
distraction: forming new internal distinctions and committing to an action
cannot happen at the same instant. An agent that is thinking is, at that
instant, not acting, and vice versa; responsiveness and reflection
alternate.

\begin{definition}[Construction and commitment]
\label{def:phases}
At each instant the agent is in one of two phases. In a \emph{construction
phase} it forms or revises internal distinctions---its disposition is
being reshaped and it commits to no response (it holds action open). In a
\emph{commitment phase} it applies a response---it commits to a single
disposition change and forms no new distinctions.
\end{definition}

\begin{axiom}[Phase exclusion]
\label{ax:exclusion}
Forming a new distinction and committing to a response draw on the same
instantaneous capacity and cannot both occur at one instant: at each
instant the agent is in exactly one phase.
\end{axiom}

\begin{theorem}[Construction--commitment alternation]
\label{thm:alternation}
Under \Cref{ax:exclusion}, an agent that both forms new distinctions and
commits responses must alternate between the two phases: there is no
instant at which it does both, so the two activities occupy disjoint
intervals in time. In particular, during a construction phase the agent
commits no responses---it is, to an outside party, unresponsive---though it
is fully active internally.
\end{theorem}

\begin{proof}
By \Cref{ax:exclusion} the instantaneous phase is single-valued, so the
sets of construction instants and commitment instants are disjoint. An
agent that does both over time must therefore partition its time into
construction intervals and commitment intervals with no overlap, i.e.
alternate. During a construction interval, by \Cref{def:phases} no
response is committed, so no external response is emitted, although the
disposition is being actively reshaped.
\end{proof}

\begin{remark}[A second, independent distraction mechanism]
\label{rem:second-distraction}
\Cref{thm:alternation} gives a source of apparent inattention distinct
from the attention-budget threshold of \Cref{sec:attention}. Even an agent
with ample budget and a relevant interaction before it will not respond
\emph{while it is in a construction phase}: it is occupied reshaping
itself and holds action open until it re-enters a commitment phase. To an
observer this is indistinguishable from being ignored, yet the agent is at
its most active. Realistic unresponsiveness therefore has two forced
origins---the budget threshold ($p^\star$) and the construction phase---and
a convincing character exhibits both.
\end{remark}

\begin{corollary}[Responsiveness is intermittent by necessity]
An agent that ever forms new distinctions is unresponsive on a positive
fraction of instants, regardless of its attention budget or the relevance
of what arrives. Continuous, immediate responsiveness is possible only for
an agent that never reflects.
\end{corollary}

\begin{proof}
Immediate from \Cref{thm:alternation}: construction intervals carry no
committed responses, and any agent that forms distinctions has construction
intervals of positive total measure.
\end{proof}

% =====================================================================
\section{Numerical study}
\label{sec:validation}
% =====================================================================

The results are proved from the axioms; we exhibit them on explicit agents
to display their shape and check the derivations directly.

\subsection{Method}

An agent is realised by: a disposition space $\Beh=[0,1]^n$; a strongly
convex drive $\drive$ with known minimiser $\Purp$ and constant $\lam$; a
finite response repertoire $\Resp$ with costs $\ge\floorc$; a consistency
graph $H$ with transports $\theta_{ij}$ giving a computable coherence
margin $\coh(x)$; and an attention budget $\att$. For streams of arriving
interactions we compute, per interaction, the best admissible response
(purpose gain and coherence test), then apply the threshold policy under
the budget. We measure: the realised response count against
$\lfloor\att/\floorc\rfloor$; the selected set against the gain-density
threshold; the pre-commitment coherence test against the actual
post-response margin; the committed-count monotonicity; and the fraction
of instants in construction versus commitment phases.

\subsection{Predictions checked}

\begin{enumerate}[label=\textbf{(V\arabic*)},leftmargin=*]
\item \textbf{Two-factor relevance.} An interaction is answered iff it has
a response clearing both bars; purpose-only and coherence-only responses
are declined (\Cref{thm:conjunction}).
\item \textbf{Inattention forced.} The committed response count never
exceeds $\lfloor\att/\floorc\rfloor$, and excess relevant interactions go
unanswered (\Cref{thm:inattention}).
\item \textbf{Threshold policy.} The value-maximising selection matches the
gain-density threshold with a single price $p^\star$; $p^\star$ rises as
$\att$ falls or arrivals grow (\Cref{thm:threshold}).
\item \textbf{Token boundary.} With small residual budget, only floor-cost
acknowledgements are committed (\Cref{cor:token}).
\item \textbf{Coherence pre-test.} The advance holonomy test predicts the
actual post-response margin sign with no error; tempting-but-destructive
responses are refused before commitment
(\Cref{thm:cohtest,cor:refuse}).
\item \textbf{Irreversibility.} Committed count is strictly increasing;
repeated encounters at equal disposition have distinct states
(\Cref{thm:irreversible}).
\item \textbf{Alternation.} Construction and commitment instants are
disjoint; response emission is zero throughout every construction interval
(\Cref{thm:alternation}).
\end{enumerate}

We report that on a sweep of randomly generated agents (varying dimension,
drive curvature $\lam$, repertoire size and costs, coherence-graph
density, and attention budget), together with hand-built witnesses for the
conjunction and coherence-refusal cases, every prediction holds: relevance
is the conjunction and neither factor alone suffices; the response count
respects the budget bound; the optimal selection is the gain-density
threshold with a single rising price; the advance coherence test matches
the realised margin exactly; the committed count is strictly monotone; and
no response is emitted during construction. (The suite is deterministic
given its inputs; figures and a reproducible harness accompany the source.)

% =====================================================================
\section{Discussion}
\label{sec:discussion}
% =====================================================================

\subsection{What has and has not been established}

We have shown that, for any agent with a purpose-directed drive, a
closed coherence condition, a floor on response cost, and a bounded
attention budget, \emph{relevance} of an interaction is the conjunction of
a purpose gain and a coherence margin; that bounded attention forces
inattention and makes the optimal policy a single-price threshold on gain
density; that coherence is testable in advance, so an agent can refuse a
tempting but self-breaking response; that response is floored and
irreversible, so encounters accumulate and no two are with the same agent;
and that construction and commitment alternate, giving a second forced
source of apparent unresponsiveness. These are theorems from the stated
axioms.

We have \emph{not} modelled how the drive itself changes as history
accumulates (the drive is fixed within an episode here), nor how an agent
chooses \emph{which} distinctions to construct during a construction phase;
both are left open. And as in the companion account, the strong-convexity
of the drive is a modelling choice that presumes a single purpose; agents
torn between purposes are out of scope.

\subsection{The two forced faces of distraction}

The paper isolates two independent reasons a well-modelled agent appears to
ignore what is put to it, neither of which is a defect. First, the
attention price $p^\star$: when the agent is busy or besieged, an
interaction below the price is correctly declined. Second, phase
alternation: while constructing, the agent commits nothing, however
relevant the interaction. A convincing embedded character exhibits both,
and a designer has, in $p^\star$ and the construction/commitment schedule,
two readable dials for how present or preoccupied a character is at any
moment.

\subsection{Relation to the account of purpose}

This paper takes purpose and coherence as given standing quantities and
asks what a relevant response to an incoming interaction is and costs. The
companion account derives where those standing quantities come from---why
an agent has an identity and a purpose at all. Together they describe an
embedded character: one paper says what the character \emph{is} and is
\emph{for}; this one says how it \emph{answers the world} in a way faithful
to what it is for.

\subsection{Closing}

An interaction matters to an agent when a response to it both moves the
agent toward its purpose and leaves the agent coherent; and because
attention is finite and reflection excludes action, the agent answers only
some of what matters and none of what does not. The result is an agent that
is present without being available---selective, occasionally
unresponsive, accumulating, and its own. None of this needs to be sold: it
is a set of theorems about finite agents responding to a world, and the
readings of them as attention and character are offered only as
orientation.

\vspace{0.6em}
\hrule
\vspace{0.4em}
\noindent\emph{Relevance is purpose gain and coherence together; attention
is priced; response is a one-way deposit; reflection and action take
turns.}

\bibliographystyle{plain}
\bibliography{references}

\end{document}
